% Template LaTeX para TCC de graduação da UNIPAMPA — V2
% Copyright (c) 2026 Charles Quevedo Carpes.
% Autoria das contribuições próprias, organização e adaptações do template.
% This work may be distributed and/or modified under the conditions of the
% LaTeX Project Public License, version 1.3c or (at your option) any later version.
% License: https://www.latex-project.org/lppl/lppl-1-3c.txt (see LICENSE).
% This work has the LPPL maintenance status `author-maintained'.
% The Current Maintainer of this work is Charles Quevedo Carpes.
% Consulte AUTORIA.md para escopo, créditos e relação dos arquivos.

% =============================================================================
% DADOS DO TRABALHO - preencha apenas este arquivo
% =============================================================================

\instituicao{Universidade Federal do Pampa}
\campus{Campus (Nome do Campus)}
\curso{(Nome do Curso)}

\autor{Nome completo do autor}
\titulo{Título do trabalho}
\subtitulo{Subtítulo, se houver}

\local{Cidade}
\data{Ano}

% Titulação pretendida e campos acadêmicos
\grau{Licenciado em Matemática}
\areaconcentracao{Nome da área, se houver}
\linhapesquisa{}

% Texto da natureza do trabalho. Mantê-lo em espaço simples.
\preambulo{Trabalho de Conclusão de Curso apresentado ao Curso de (Nome do curso) da Universidade Federal do Pampa, como requisito parcial para obtenção do Título de Licenciado (ou Bacharel) em (Área do curso).}

\orientador{Prof. Dr. Nome completo do orientador}
% Deixe vazio quando não houver coorientador.
\coorientador{Nome do Coorientador}

\datadefesa{dia de mês de ano}

% Cada membro: nome e titulação; função; instituição.
\membrobanca{Prof. Dr. Nome completo do orientador}{Orientador}{Universidade Federal do Pampa}
\membrobanca{Profa. Dra. Nome do segundo membro}{Coorientador - Opcional}{Universidade Federal do Pampa}
\membrobanca{Prof. Dr. Nome do terceiro membro}{Instituição do membro}

% Logo é opcional e não é usado por padrão.
% \logo{figuras/logo-unipampa.pdf}

% Ficha catalográfica gerada no GURI. Enquanto o arquivo não existir, o template
% cria uma página de reserva. A ficha é excluída da contagem de páginas.
\arquivoFichaCatalografica{documentos/ficha-catalografica.pdf}

% Para o depósito final, é possível substituir a folha de aprovação editável por
% uma página assinada em PDF. Deixe comentado enquanto estiver redigindo.
% \arquivoFolhaAprovacao{documentos/folha-aprovacao-assinada.pdf}
